\documentclass[journal]{IEEEtran}

\usepackage{amsmath,amssymb,amsthm}
\usepackage{bm}
\usepackage{graphicx}
\usepackage{booktabs}
\usepackage{cleveref}

\newtheorem{proposition}{Proposition}

\newcommand{\rr}{\bm{r}}
\newcommand{\xx}{\bm{x}}
\newcommand{\hh}{\bm{h}}
\newcommand{\Gtil}{\widetilde{\mathbf{G}}}
\newcommand{\Sab}{S_{\mathrm{ab}}}
\newcommand{\SARwb}{\mathrm{SAR}_{\mathrm{wb}}}
\newcommand{\SAR}{\mathrm{SAR}}
\newcommand{\APD}{\mathrm{APD}}
\newcommand{\Pabs}{P_{\mathrm{abs}}}
\newcommand{\Q}{\mathbf{Q}}
\newcommand{\Aeff}{A_{\mathrm{eff}}}

\begin{document}

\title{From whole-body SAR to spatial compliance: \\
  $\APD_{4}$ almost always binds first in mmWave MIMO}
\author{Opus (draft for direction 4)}
\maketitle

\begin{abstract}
At mmWave, ICNIRP replaces the whole-body SAR constraint by a
$4\,\mathrm{cm}^2$-averaged absorbed power density (APD) constraint
applied everywhere on the body surface. Using the monograph's coherent
exposure operator we cast both as quadratic forms in the precoder,
derive a scalar crossover condition based on the ``effective
illumination area'' $\Aeff$, and verify the implied scaling with an
AEGIS simulation of an $N$-element MIMO hotspot on a flat skin panel.
The critical area for an adult is $\Aeff^\star=0.28\,\mathrm{m}^2$.
Any beam-shaped illumination that lights less than this surface --
i.e.\ effectively any mmWave link -- puts the system in the
APD-dominated regime, and coherent MIMO beamforming moves it deeper
in by a factor that scales linearly with the array size once the
hotspot exceeds the $4\,\mathrm{cm}^2$ patch.
\end{abstract}

%======================================================================
\section{Problem}\label{sec:problem}
%======================================================================

Consider an $M$-antenna UE or BS transmitting with precoder
$\xx\in\mathbb{C}^M$ ($\lVert\xx\rVert^2\le P$) over a channel
$\hh\in\mathbb{C}^M$. The body is a surface $\Sigma\subset\mathbb{R}^3$
with outward normal $\hat{\bm n}(\rr)$, mass $m$, and area
$|\Sigma|$.

From the monograph (Def.~5.1, Thm.~4.1), the absorbed power density at
$\rr\in\Sigma$ is
\begin{equation}\label{eq:Sab}
  \Sab(\rr) \;=\; \bigl\lVert\Gtil(\rr)\,\xx\bigr\rVert^2
  \;=\; \xx^H\bigl(\Gtil(\rr)^H\Gtil(\rr)\bigr)\xx ,
\end{equation}
with $\Gtil(\rr)\in\mathbb{C}^{3\times M}$ the Fresnel-filtered,
depth-coupled body-surface channel. Integrating over $\Sigma$ yields
the global exposure operator $\Q=\int_\Sigma\Gtil^H\Gtil\,\mathrm{d} A$
and $\Pabs=\xx^H\Q\xx$.

Public-exposure limits switch regime at $6$\,GHz:
\begin{align}
  \text{(below 6\,GHz, bulk)}\quad&
  \SARwb \;=\; \Pabs/m \;\le\; L_{\SARwb}^\star/m, \label{eq:SARwb}\\
  \text{(above 6\,GHz, surface)}\quad&
  \APD_4(\rr_0)\;\le\; L_{\APD} \;\;\forall\rr_0\!\in\!\Sigma,\label{eq:APDpt}
\end{align}
with $L_{\SARwb}^\star=0.08\,\mathrm{W/kg}$, $L_{\APD}=20\,\mathrm{W/m^2}$,
$A_p=4\,\mathrm{cm}^2$, and
\begin{equation}\label{eq:APDdef}
  \APD_4(\rr_0) \;=\; \frac{1}{A_p}\!\int_{\mathcal{P}(\rr_0,A_p)}\!\Sab(\rr)\,\mathrm{d} A.
\end{equation}
FR2 (24--71\,GHz) and 6G candidate bands lie above the threshold:
\eqref{eq:SARwb} is no longer the governing constraint there;
\eqref{eq:APDpt} is. The balance between the two was examined at
sub-6\,GHz by Hochwald \emph{et al.}~\cite{Hochwald2014} and
Ying~\emph{et al.}~\cite{Ying2015} through a single quadratic
``SAR matrix'' $\mathbf S_V$; above 6\,GHz no such scalar exists.

%======================================================================
\section{Local exposure operator}\label{sec:local-Q}
%======================================================================

Substituting \eqref{eq:Sab} into \eqref{eq:APDdef},
\begin{equation}\label{eq:Qloc}
  \APD_4(\rr_0) \;=\; \xx^H\,\Q_{\mathrm{loc}}(\rr_0)\,\xx,
  \;\;
  \Q_{\mathrm{loc}}(\rr_0) = \frac{1}{A_p}\!\int_{\mathcal{P}(\rr_0,A_p)}\!\Gtil^H\Gtil\,\mathrm{d} A .
\end{equation}
Each $\Q_{\mathrm{loc}}(\rr_0)\in\mathbb{C}^{M\times M}$ is Hermitian
positive-semidefinite. The family
$\{\Q_{\mathrm{loc}}(\rr_0)\}_{\rr_0\in\Sigma}$ replaces the single
$\Q$ in ECBF (monograph Sec.~8).

On a discretised mesh with $M_\triangle$ triangles, centroids
$\{\rr_i\}$, areas $\{a_i\}$, the ICNIRP spatial averaging is a
row-stochastic sparse matrix $W\in\mathbb{R}^{M_\triangle\times M_\triangle}_{\ge 0}$
with row $i$ supported on the minimal ball around $\rr_i$ whose
integrated area reaches $A_p$. Then
\begin{equation}\label{eq:QlocD}
  \Q_{\mathrm{loc}}(\rr_i)
  \;=\; \sum_{j} W_{ij}\,\Gtil(\rr_j)^H\Gtil(\rr_j),
\end{equation}
and the full family is a linear reweighting of $M_\triangle$ rank-$\!\le\! 3$
matrices. For a fixed $\xx$, the worst-case patch is computed directly
from the per-triangle $\bm{s}_i=\Sab(\rr_i)$ as $\max_i (W\bm{s})_i$,
costing one sparse mat-vec.

The routine building $W$ exists in AEGIS as
\texttt{aegis.geometry.averaging.\allowbreak precompute\_averaging\_matrix};
it returns a Numba-accelerated CSR matrix that can be cached across
precoders for the same mesh.

%======================================================================
\section{The joint-constrained problem}\label{sec:jointP}
%======================================================================

Writing $L_{\SAR}\!:=\!m\,L_{\SARwb}^\star$ (total body
absorption budget), the rate-maximising precoder with both regulatory
regimes active is
\begin{equation}\label{eq:P}
  \boxed{\;
  \begin{aligned}
    \max_{\xx\in\mathbb{C}^M}\;& |\hh^H\xx|^2 \\
    \text{s.t.}\;& \xx^H\Q\xx \le L_{\SAR} &&\text{(whole-body SAR)}\\
    & \xx^H\Q_{\mathrm{loc}}(\rr_0)\xx \le L_{\APD} \;\;\forall\rr_0\!\in\!\Sigma && \text{(local APD)} \\
    & \lVert\xx\rVert^2 \le P &&\text{(transmit power)}.
  \end{aligned}\;}
\end{equation}
On a discretised body this is a non-convex QCQP with
$1+M_\triangle+1$ quadratic constraints. The Hochwald--Ying closed
form~\cite{Hochwald2014,Ying2015} and the monograph's Sec.~8 handle
the special case $M_\triangle\!=\!0$; with $M_\triangle\!\gg\! 1$
S-procedure duality is no longer guaranteed tight in general, but
for a single dominant hotspot only the tightest
$\Q_{\mathrm{loc}}(\rr_\star)$ is active at the optimum, reducing
\eqref{eq:P} to a two-constraint QCQP with the same closed-form
solution $\xx^\star\!\propto\!(\lambda_1\Q_{\mathrm{loc}}(\rr_\star)+\lambda_2\Q+\nu\mathbf I)^{-1}\hh^*$.
Identifying $\rr_\star$ is the real problem; it is a non-convex search
over $\Sigma$ that the ECBF loop has to iterate with.

The structural takeaway: $\Q$ averages over the body and is blind to
hotspots; $\Q_{\mathrm{loc}}$ localises. Whichever of the
$M_\triangle\!+\!1$ exposure constraints is tightest at unconstrained
MRT sets the binding one.

%======================================================================
\section{When does the local limit bind?}\label{sec:crossover}
%======================================================================

The ratio of the two limits controls the crossover. Define the
precoder's \emph{coupling efficiency}
\begin{equation}\label{eq:kappa}
  \kappa(\xx) \;:=\; \frac{\max_{\rr_0\in\Sigma}\APD_4(\rr_0)}{\Pabs}
  \;=\; \frac{\lambda_{\max}\!\bigl\{\Q_{\mathrm{loc}}(\rr_0)\bigr\}_{\rr_0}}{\xx^H\Q\xx/\lVert\xx\rVert^2}
  \quad[\mathrm{m}^{-2}]
\end{equation}
and the \emph{effective illumination area}
$\Aeff(\xx):= 1/\kappa(\xx)$. For a uniformly illuminated
surface of area $A$ embedded in a larger shadowed body, $\kappa=1/A$
exactly, so $\Aeff=A$; this justifies the name. For a focused
hotspot of area $A_h\le A_p$ on a uniformly dim background,
$\Aeff\!\to\! A_p$ (averaging floor); for a hotspot broader than
$A_p$, $\Aeff\!\to\! A_h$.

\begin{proposition}[Crossover]\label{prop:cross}
The local APD constraint binds strictly before whole-body SAR iff
\begin{equation}\label{eq:crossover}
  \Aeff(\xx) \;<\; \Aeff^\star
  \;:=\; \frac{L_{\SAR}}{L_{\APD}}
  \;=\; \frac{m\,L_{\SARwb}^\star}{L_{\APD}}.
\end{equation}
For the ICNIRP public-exposure reference: adult
$(m\!=\!70,|\Sigma|\!=\!1.8\,\mathrm m^2)$, $\Aeff^\star\!=\!0.28\,\mathrm m^2$
($15.6\%$ of $|\Sigma|$); infant $(m\!=\!8,|\Sigma|\!=\!0.4)$,
$\Aeff^\star\!=\!0.032\,\mathrm m^2$ ($8.0\%$); large adult
$(m\!=\!100,|\Sigma|\!=\!2.0)$, $\Aeff^\star\!=\!0.40\,\mathrm m^2$
($20\%$).
\end{proposition}

\begin{proof}
At SAR saturation $\Pabs\!=\!L_{\SAR}$, the worst-case APD is
$\kappa L_{\SAR}$. APD binds before SAR iff
$\kappa L_{\SAR}\!>\!L_{\APD}$, i.e.\ $\kappa\!>\!L_{\APD}/L_{\SAR}$,
i.e.\ $\Aeff\!<\!L_{\SAR}/L_{\APD}$.
\end{proof}

\Cref{prop:cross} recasts the regulatory regime question as a purely
geometric one: a beam that lights up more than a quarter of an adult's
skin is SAR-limited; anything tighter than that is APD-limited.
In practice all mmWave scenarios -- a handset held at 5--30\,cm, a
5G CPE on a desk, a BS steering a pencil beam at a user -- illuminate
body areas well below $0.28\,\mathrm m^2$; the smaller the focus, the
further into the APD-dominated regime the system lies.

\subsection{Scaling of $\Aeff$ with MIMO}

$N$ antennas radiating co-polarised plane waves from a cone of
half-angle $\alpha$ toward the body with MRT phasing produce a
coherent hotspot at the cone axis. The monograph
(Thm.~4.1 and Table~4) gives a peak coherent gain of $N$ relative to
the body-average $\Sab$; combined with the hotspot coherence length
$\delta_{\mathrm{coh}}\!\sim\!\lambda/(2\pi\sigma_\Omega)$ (Sec.~7.2)
this gives two regimes for the $4\,\mathrm{cm}^2$-averaged peak:

\noindent\emph{Narrow cone} ($\sigma_\Omega$ small,
$\delta_{\mathrm{coh}}^2\!\gtrsim\!A_p$): the hotspot is broader than the
averaging patch, so averaging does not dilute the peak. Then
$\Aeff\!\sim\!|\Sigma_{\mathrm{lit}}|/N$, where
$|\Sigma_{\mathrm{lit}}|$ is the lit body area; $\Aeff$ shrinks as
$1/N$.

\noindent\emph{Wide cone} ($\delta_{\mathrm{coh}}^2\!<\!A_p$): the
hotspot fits inside the $4\,\mathrm{cm}^2$ patch; averaging dilutes
the peak by the fill factor $\delta_{\mathrm{coh}}^2\pi/A_p$.
$\Aeff$ saturates near $|\Sigma_{\mathrm{lit}}|A_p/(N\delta_{\mathrm{coh}}^2\pi)$
in the best case for compliance but never above $A_p/N$ coherent
multiplication times $\pi\delta_{\mathrm{coh}}^2$ in absolute value.

\subsection{AEGIS numerical verification}

We verify the scaling on a flat $30\!\times\!30$\,cm skin panel
($28\,800$ isoceles triangles at $2.5$\,mm pitch, \texttt{SKIN\_28GHZ}
dielectric) illuminated by $N$ plane waves from directions uniform in
a cone of half-angle $\alpha$ around $-\hat{\bm z}$, co-polarised
along $\hat{\bm x}$, co-phased at the panel centre via
$h_j\!=\!1$. Code is in
\verb|JSAC/direction_4_sim_opus/hotspot_scaling.py|; averaging via
\texttt{precompute\_averaging\_matrix}; three seeds averaged per
$(N,\alpha)$.

\Cref{tab:eta4} reports the panel-level ratios $\eta_4^{\mathrm{peak}}\!=\!\Sab^{\mathrm{peak}}/\bar\Sab_{\mathrm{panel}}$
(no averaging) and $\eta_4\!=\!\APD_4^{\mathrm{peak}}/\bar\Sab_{\mathrm{panel}}$
(ICNIRP averaging). Since the panel is the only absorbing region in
the scenario, $\kappa = \eta_4/A_{\mathrm{panel}}$ where
$A_{\mathrm{panel}}\!=\!0.09\,\mathrm{m}^2$; the last column shows
$\Aeff\!=\!1/\kappa$ in cm$^2$.

\begin{table}[!t]
\centering
\caption{AEGIS simulation of $N$-antenna MIMO hotspot on a flat
$30\!\times\!30$\,cm skin panel at $28$\,GHz, MRT phasing at centre,
three seeds per cell. For each $\alpha$, left column is
$\eta_4^{\mathrm{peak}}$ (pointwise $\Sab$ peak / panel average), right
column is $\eta_4$ (ICNIRP $4\,\mathrm{cm}^2$ peak / panel average).
Last column is $\Aeff=A_{\mathrm{panel}}/\eta_4$ for $\alpha=10^\circ$;
values below $\Aeff^\star\!=\!2\,800\,\mathrm{cm}^2$ (\cref{prop:cross})
are APD-limited. Every entry satisfies $\Aeff\!<\!\Aeff^\star$: APD
binds everywhere.}
\label{tab:eta4}
\begin{tabular}{r cc cc cc c}
\toprule
 & \multicolumn{2}{c}{$\alpha\!=\!10^\circ$}
 & \multicolumn{2}{c}{$30^\circ$}
 & \multicolumn{2}{c}{$60^\circ$}
 & $\Aeff$ \\
\cmidrule(lr){2-3}\cmidrule(lr){4-5}\cmidrule(lr){6-7}
$N$ & $\eta^{\mathrm{pk}}$ & $\eta_4$ & $\eta^{\mathrm{pk}}$ & $\eta_4$
    & $\eta^{\mathrm{pk}}$ & $\eta_4$ & [cm$^2$] @ $10^\circ$\\
\midrule
  1 &  1.0 &  1.0 &  1.0 &  1.0 &  1.0 &  1.0 &    900 \\
  2 &  2.1 &  2.1 &  2.0 &  1.6 &  1.9 &  1.3 &    439 \\
  4 &  3.5 &  3.0 &  3.3 &  1.8 &  2.8 &  1.2 &    299 \\
  8 &  7.7 &  6.7 &  7.5 &  3.1 &  6.2 &  1.6 &    134 \\
 16 & 12.7 & 10.8 & 14.9 &  5.1 & 11.4 &  2.6 &   \phantom{0}84 \\
 32 & 22.1 & 19.0 & 28.8 & 10.0 & 23.0 &  4.1 &   \phantom{0}47 \\
 64 & 36.1 & 30.8 & 53.3 & 17.9 & 43.7 &  7.4 &   \phantom{0}29 \\
\bottomrule
\end{tabular}
\end{table}

Three qualitative observations match \cref{prop:cross} and the
scaling argument.

\emph{(i)} At $\alpha\!=\!10^\circ$ the $4\,\mathrm{cm}^2$ averaging
has almost no effect ($\eta_4\!\approx\!0.85\,\eta^{\mathrm{pk}}$): the
hotspot is broader than the patch, so $\eta_4$ tracks $N$
($\eta_4\!\approx\!0.48\,N$, slope matching the $N$-scaling of
Table~4 in the monograph).

\emph{(ii)} At $\alpha\!=\!60^\circ$ the hotspot is sub-wavelength and
much smaller than $A_p$, so the $4\,\mathrm{cm}^2$ averaging dilutes
the peak roughly by $A_h/A_p$; $\eta_4$ grows sublinearly
($\sim N^{0.5}$) and saturates at a much lower value. Not shown in
the table: the pointwise $\eta^{\mathrm{pk}}$ still grows linearly with
$N$; the crop is an averaging artefact, not a physical protection.

\emph{(iii)} The last column shows that the single-antenna
($N\!=\!1$) case already has $\Aeff\!=\!900\,\mathrm{cm}^2$, about
$3.1\times$ below $\Aeff^\star\!=\!2\,800\,\mathrm{cm}^2$, so
\emph{even without any MIMO gain} the APD limit is already tighter
than SAR for this geometry -- concentrating the illumination onto
$9\%$ of the body surface is enough. Each doubling of $N$ then buys
another factor of 2 (narrow cone) or less (wide cone) of
concentration.

%======================================================================
\section{Implications}\label{sec:implications}
%======================================================================

\emph{(a) Whole-body SAR is not the design driver above 6\,GHz.} Any
beam-shaped illumination with $\Aeff\!<\!0.28\,\mathrm{m}^2$ is
APD-limited, and all realistic mmWave scenarios satisfy that.
Papers using a single scalar $\Q$ (or $\mathbf S_V$) to bound
exposure are optimistic in this regime:
they tolerate precoders with feasible $\xx^H\Q\xx$ but infeasible
$\xx^H\Q_{\mathrm{loc}}(\rr_\star)\xx$.

\emph{(b) ECBF must track hotspots.} The single-constraint ECBF
closed form is easy to state but spatially blind. Solving~\eqref{eq:P}
demands alternating (\emph{i}) hotspot identification
$\rr_\star(\xx)=\arg\max_{\rr_0}\xx^H\Q_{\mathrm{loc}}(\rr_0)\xx$
and (\emph{ii}) two-constraint ECBF with
$\{\Q_{\mathrm{loc}}(\rr_\star),\Q\}$. The monograph's
$\lambda_1\Q_{\mathrm{loc}}+\lambda_2\Q+\nu\mathbf I$ structure still
applies per iteration.

\emph{(c) Finite covering of $\Sigma$ is tractable.} Because
$\Q_{\mathrm{loc}}(\rr_0)$ is Lipschitz in $\rr_0$ over scales
$\gtrsim\!\sqrt{A_p}/2\!\approx\!1\,\mathrm{cm}$ for mmWave field
channels, only $\sim\!|\Sigma|/A_p\!\approx\!4\,500$ patch centres
cover an adult body; fewer in practice because only hotspot
candidates matter. AEGIS already computes the per-triangle $\Sab$ map,
so the worst patch is one sparse mat-vec away: the cost of promoting
ECBF from global to spatial is modest.

\emph{(d) Zhou \emph{et al.}~\cite{Zhou2026}'s thermal relaxation sits
on top.} Their long-term bioheat constraint replaces the rigid
$\APD_4\!\le\!L_{\APD}$ by a queue-stability condition. Whichever
constraint is in effect (instantaneous in this note, time-averaged
there) is still a family indexed by $\rr_0\in\Sigma$, not a single
matrix. The spatial-covering framework developed here applies
unchanged and is a prerequisite to a physically-correct thermal ECBF.

%======================================================================
\bibliographystyle{IEEEtran}
\begin{thebibliography}{9}
\bibitem{ICNIRP2020} ICNIRP, ``Guidelines for limiting exposure to
  electromagnetic fields (100\,kHz--300\,GHz),'' \emph{Health Phys.},
  vol.~118, no.~5, pp.~483--524, 2020.
\bibitem{IEEE95_2019} IEEE Std C95.1-2019, ``IEEE Standard for Safety
  Levels with Respect to Human Exposure to Electric, Magnetic, and
  Electromagnetic Fields, 0\,Hz to 300\,GHz,'' 2019.
\bibitem{Hochwald2014} B.~Hochwald, D.~Love, S.~Yan, P.~Fay, and
  J.~Jin, ``Incorporating specific absorption rate constraints into
  wireless signal design,'' \emph{IEEE Commun. Mag.}, vol.~52, no.~9,
  pp.~126--133, 2014.
\bibitem{Ying2015} D.~Ying, D.~Love, and B.~Hochwald, ``Closed-loop
  precoding and capacity analysis for multiple-antenna wireless
  systems with user radiation exposure constraints,'' \emph{IEEE
  Trans.\ Wireless Commun.}, vol.~14, no.~10, pp.~5859--5870, 2015.
\bibitem{Zhou2026} Z.~Zhou \emph{et al.}, ``Exposure-aware beamforming
  for mmWave systems: from EM theory to thermal compliance,''
  arXiv:2601.19587, 2026.
\end{thebibliography}

\end{document}
